\section{Critical Evaluation, Conclusion, and Future Work}\label{sec:chapter-8}

\subsection{Evaluation against objectives}

The source and recorded checks support O1 through implemented account, onboarding, check-in, and persistence pathways, although complete release testing remains open. O2 is supported by a working rules-and-semantic pipeline, recorded score components, and baseline evaluation. O3 has specific evidence for declared exclusions, including short or empty candidate handling. O4 is supported by implemented device encryption and documented native/\allowbreak{}backend round trips. O5 is supported by the administration and content modules. O6 is partially achieved because technical evaluation exists, while representative usability and independent suitability assessment do not.

The strongest engineering result is consistency across boundaries. Declared comfort concerns become explicit constraints, filtered candidates remain excluded during selection, and encrypted journals do not quietly re-enter server recommendation context. The artefact also supports content administration without requiring every content change to become a mobile code change. These outcomes demonstrate useful integration. They do not prove that all content is appropriate or that every deployment path is equally reliable.

\subsection{Threats to validity}

The benchmark is small and authored within the project, so selection and label bias are plausible. A fixed catalogue limits coverage, and preserved labels may reflect expectations inconsistent with newly declared exclusions. Cold-start scenarios cannot assess long-term collaborative behaviour or rotation quality. Repeating deterministic rankers primarily supports timing; it does not enlarge the independent sample. Rules and text baselines differ in eligibility behaviour, so measured improvements cannot be assigned solely to semantic weighting.

Other evidence boundaries concern environment and time. Simulator validation does not establish physical-device signing, interruption handling, reinstall recovery, or Android behaviour. Historical documentation may lag current source. Warm ranking latency excludes network and interface delay. Cryptographic implementation evidence does not replace an independent security review, and current row migration does not erase earlier backups. These limitations should constrain both academic conclusions and public product claims.

\subsection{Usability, maintenance, and practical impact}

The application's practical value lies in reducing the steps between noticing a state and finding an available activity. That value remains plausible rather than measured. A user study should examine whether check-in questions are understandable, whether comfort choices accurately describe what people wish to avoid, and whether recommendation reasons help users make decisions. It should also observe skipped questions and declined practices, because these behaviours may reveal friction that a favourable satisfaction score alone would miss.

Maintenance is equally relevant to future impact. A published exercise can change wording, media, or metadata without changing the ranking code. Such updates may alter semantic similarity or eligibility, making catalogue versioning and regression checks valuable. Journal envelope versions and native cryptography dependencies also require coordinated compatibility. The project should retain reproducible evaluation snapshots while documenting when their catalogue or engine no longer matches a release. This would improve the usefulness of both favourable and unfavourable historical results.

\subsection{Comparison with the reviewed knowledge base}

The artefact applies the combination principle described by \citet{burke2002} through an eligibility gate and a conditional weighted ranker. Its contribution lies in the integration of those mechanisms with a daily mobile workflow, rather than a new general-purpose recommendation algorithm. It uses the pretrained semantic representation described by \citet{reimers2019}, without claiming domain-specific model training. The saved comparison shows a 4.17 percentage-point Precision@4 improvement over rules, but that improvement accompanies a higher negative-label selection rate and does not establish a better experience for users.

The NarraGive work described by \citet{slade2024} motivates attention to several evaluation dimensions rather than a single accuracy score. Holistic Mind reports relevance, ordering, violations, coverage, diversity, and timing. However, its authored, interaction-free fixtures provide a narrower evidence base than an evaluation of actual narrative use. The project's collaborative component exists in implementation but remains untested by the saved cold-start benchmark. Independent access through Explore provides a useful design route outside personalised ranking, although its effect on feedback bias has not been measured.

The mixed evidence reviewed by \citet{matthews2025} also limits interpretation: personalisation does not automatically imply additional mental-health benefit. Holistic Mind demonstrates implemented selection support and a privacy boundary, not a proven therapeutic intervention. Relative to a recommender-platform approach such as \citet{tapuria2024}, it concentrates on selection within a managed exercise catalogue and connects that selection to practice feedback and encrypted reflection. These comparisons identify the artefact's scope and unresolved questions without claiming superiority over systems evaluated on different tasks.

\subsection{Answers to research questions}

RQ1 is answered provisionally: on the journal-free authored benchmark, hybrid ranking improved precision and NDCG over the included alternatives, but also returned more negative-labelled items than rules or semantic-only ranking. RQ2 is supported for declared restrictions by the comfort suite and targeted regression cases; undeclared restrictions remain unresolved. RQ3 is supported at implementation level: encrypted journals persist through the API while the current recommendation route sends no journal text. The resulting loss of reflective context is an intentional design trade-off rather than a missing server decryption feature.

\subsection{Future research and privacy work}

The next priority is independent review of exercise metadata and relevance labels, followed by weight selection on development scenarios and evaluation on unseen cases. A longitudinal dataset should assess feedback adaptation, collaborative activation, and repeated recommendation quality. Physical iOS and Android tests should cover interruptions, secure-storage behaviour, recovery on another device, media, and account deletion. A consented usability study should measure task completion and perceived control, while separate accessibility checks examine contrast, assistive technology, and text scaling.

Further privacy work should examine key rotation, compromised recovery material, background locking, and retention of legacy backups. Any journal-aware personalisation should begin with a distinct local-processing design and its own performance and privacy evaluation. Holistic Mind demonstrates that an integrated wellness application can combine a practical daily workflow, constrained recommendation logic, and device-side journal privacy. Its saved results justify continued development, while the reported violations and evidence gaps identify exactly where stronger validation is still required.

\subsection{Reflection on development decisions}

The strongest continuity across the earlier submissions is the commitment to a supportive, non-diagnostic interaction. The June review's concern about sensitive onboarding shaped later choices to collect short structured answers rather than a personal trauma narrative. Its question about explainable recommendations is addressed through returned reasons and score components. Its question about evaluation beyond random suggestions is addressed through four baseline approaches. The final system therefore develops several early research questions into concrete mechanisms and evidence.

Replacing Firebase was the most substantial architectural change reported in the final progress review. It introduced redevelopment of authentication, data access, and administration, but enabled direct control over relational records and recommendation coordination. The decision is defensible for the resulting artefact, although the supplied material does not document a formal total-cost or security comparison with Firebase. The report should explain the observed benefits without treating a custom backend as inherently superior. Greater control also transfers operational and security responsibilities to the project.

The journal redesign illustrates a second important learning point. An apparently useful recommendation input can conflict with the product's privacy intention. Removing journal context reduced the available semantic evidence and slightly lowered saved NDCG compared with the historical benchmark. The project addressed that trade-off by publishing a journal-free evaluation instead of continuing to promote scores from an input path the app no longer uses. This is stronger research practice than retaining a favourable but mismatched historical result.

Native debugging also changed the understanding of completion. The required module and keychain signing problems showed that an exported JavaScript bundle is not equivalent to a working installed app. The stale recommender container showed that source version and running version can diverge. The October sign-in issue showed that provider configuration is insufficient when the application server is stopped. Each problem crossed a boundary between software layers, supporting the final review's emphasis on systems thinking and maintainability.

\subsection{Practical impact and commercial scope}

The original proposal's freemium idea describes one possible business model, but the current academic contribution does not depend on subscriptions. Introducing paid access would require payment and entitlement handling, clear content boundaries, and a decision about which practices remain available without charge. It would also introduce new usability and retention questions. None of those questions is resolved by the existing recommendation metrics. The report therefore treats monetisation as optional future product work.

Potential impact is concentrated in a coherent routine, clearer content selection, and greater control over reflective privacy. These benefits are plausible from the design but need observation with intended users. A larger library could also create new selection burden if categories and reasons become confusing. More engagement is not automatically better: repeated use might indicate usefulness, habit, distress, or simple curiosity. The project should avoid optimising solely for completion counts and instead assess what users want from the routine.

\subsection{Prioritised completion roadmap}

The first priority is content and metadata review. An independent reviewer should examine descriptions, restrictions, suitability labels, and the three unresolved benchmark cases. The output should be a dated, versioned catalogue and documented label decisions. The second priority is physical-device and release verification, covering provider sign-in, native key storage, recovery on another device, interruptions, attachments, and deletion. These activities close important evidence gaps without introducing a new recommendation feature.

The third priority is a consented formative usability study focused on the complete daily journey and journal setup. It should measure successful task completion, the interpretation of recommendation reasons, and understanding of comfort choices and key-loss consequences. The fourth is ranking research: separate development and unseen scenarios, compare common eligibility policies where appropriate, and evaluate longitudinal feedback and rotation. Only after these foundations are stronger should the project increase collaborative influence or add another personalisation input.

Future journal-aware support should be designed as a separate local-processing capability if pursued. It must not weaken the existing encryption promise by quietly uploading decrypted content or embeddings. Key rotation, recovery-material compromise, background locking, and legacy-backup retention are also concrete privacy priorities. This roadmap preserves the original ambition while directing effort toward validation and control. Holistic Mind's contribution is a functioning, inspectable artefact and a clear account of how its design evolved; its remaining research and release work are identifiable rather than concealed.

\subsection{Conclusion}

Holistic Mind translates an early proposal into an integrated mobile application with an inspectable recommendation pipeline, structured content administration, and device-side journal encryption. The project's development history shows how practical constraints and privacy decisions reshaped the original design. The saved benchmark supports the feasibility of combining rules and semantic similarity, while its negative-labelled selections demonstrate why ranking performance alone is insufficient. The next stage should strengthen content review, device verification, and user acceptance evidence. The report therefore establishes a substantial software contribution and a credible basis for further research without claiming clinical benefit or completed studies that the evidence does not support.
